\chapter{两座城}
\chapterkana{二つの都}
\epigraphbox{“月者，阴精之宗。积而成兽，象兔。”}{——张衡《灵宪》，《后汉书·天文志》刘昭注引}

\kaishou{照片上的月亮，跟平时见过的样子不太一样。\par{}那几片浅色的洼地，边太圆了，圆得不太像是撞出来的。\par{}笔者动手量过一次，量到一半就停手了。这一章写的正是那半截。}

\section{二十九帧}

一九五九年十月七日，苏联的月球三号掠过月球背面，拍下二十九帧照片，覆盖了约七成背面的地貌。\footnote{月球三号同年十月四日发射，十月七日拍摄；同月二十七日塔斯社公布首批影像，一九六〇年苏联科学院出版《月球背面图集》，二十九帧尽收其中。}胶片在探测器上直接显影，再转成电视信号发回地面，最清楚的一帧分辨率大约一公里。在此后二十年里，人类对月球背面的全部认识都靠这二十九帧影像撑着。

在那之前，学界普遍认为月面两侧大致一样，都是明亮的高地夹着几片暗淡的月海。二十九帧传回来之后，格局就对不上了：正面那些暗色的玄武岩平原（古人按颜色叫作海）大约占三成面积，背面却不到百分之三。就算按最宽的口径估算，两侧也差着一个数量级。\footnote{按此后各次环月探测的统计，正面月海约占百分之三十一，背面约百分之一。本书取「不到百分之三」这一宽口径，免得把话说满。}

光说低地稀少，其实算不上什么疑点。小天体撞击在月面两侧的概率本来差不多，玄武岩熔岩只在月壳裂隙多的一侧往外涌，所以正面裂隙多、月海多，背面低地少，这完全符合行星演化的常规。真正让人犯疑的，是那仅有的几片低地本身的形状。背面最大的一片横跨约九百公里，最小的一片约一百八十公里，表面都比周围高地平得多。按照常理，它们应当是被熔岩填平的撞击盆地，可是把它们摆在同一张底片上比对，四条边界却勾得比周围任何天然地形都要齐整。

\section{边界的量法}

要判断一条边界是自然长出来的还是人工划出来的，最省事的办法就是把它当成一条几何曲线来量。\footnote{把曲线在不同尺度上求曲率并看它的分布，是图像处理里的常规办法，一般叫曲率尺度空间；本书只用到一维形式。}做法很简单，先在影像上沿着低地与高地的交界逐点拟合局部曲率半径，再算这一串半径的离散程度。山脊或撞击坑缘这类自然边界，曲率半径一会儿大一会儿小，离散度很高；人工划出来的边界只要按固定规矩走，局部曲率半径在各处就会相差无几。笔者不是专门搞行星地质出身，手里也没有更精细的原始成像工具，只能用这套最朴素的算子，拿同一批月球三号影像测了两组样本。目标低地边界取一百二十个测点，算出来的曲率半径标准差在百分之三以内；作为对照，笔者又在高地上找了同等规模的自然边界，用同样的点数和同样的算法量了一遍，标准差达到了百分之十九。

\begin{table}[htbp]
  \centering
  \caption{同一批影像上两种边界的曲率测量}
  \begin{tabular}{@{}lcc@{}}
    \toprule
    & 目标低地边界 & 高地自然边界\\
    \midrule
    测点数 & 120 & 120 \\
    曲率半径标准差 & 3\%以内 & 19\% \\
    \bottomrule
  \end{tabular}
  \tabsource{据月球三号二十九帧影像测量。影像分辨率约一公里，测点位置误差与像元同量级。}
\end{table}

这套量法有一个明摆着的反例得先排除掉：熔岩充填一样能把地形抹平。正面月海内部的起伏很小，光看表面平不平，与背面那几片低地拉不开差距，所以「平」本身作不了证。两边的区别出在边界走向。熔岩往哪里流全看地势，哪块低洼就往哪块漫过去，边界自然顺着等高线弯弯曲曲，形状很杂乱。背面那几片低地却根本不看地势，有几段边界甚至直接横着切过高地的隆起。一条边界接近正圆，还毫不避让原有的地形起伏，画出它的东西就不太像是月表自己长出来的地质过程。

当然，这套测量本身也很粗糙。一百二十个测点全都点在分辨率约一公里的影像上，一个像元就是一公里见方，取点时手稍微偏一点，误差就能占掉百分之几。这样的误差能抹掉一部分精度，只是还不足以把百分之三和百分之十九的差距拉平。

\section{壳层之下}

边界齐整只能说明它被人划定过，还看不出底下的构造到底是什么。要查清那几片低地下面有没有被撞击挖空过，还得靠重力场数据。二〇一一年至二〇一二年，圣杯号双星沿月球轨道编队飞行，通过连续测量两颗卫星间距的微小变化，反演出了高精度的月球重力场。\footnote{圣杯号（GRAIL）二〇一一年九月发射，二〇一二年十二月受控撞月。本书采用之后公开的重力场模型，没有处理原始跟踪数据。}

这项观测给出了一条从前没有的判据：正面几乎每一片月海下面都有质量瘤，也就是玄武岩充填加上深部壳层抬升共同造成的正重力异常。\footnote{质量瘤最早由 Muller 与 Sjogren 在一九六八年从轨道跟踪数据里算出来。后来各家复算过几次，正负号一致，数值略有出入。}它的来历很清楚：大型撞击盆地形成时会掀掉表层月壳，逼着下面的地幔往上抬，随后坑里又灌满高密度熔岩，重力场上自然就会留下正异常的印记。

背面那几片低地下面却找不到质量瘤。标准的撞击盆地剖面有固定的几件东西：外缘隆起，内部塌陷，坑底下面的月壳变薄，正中常常还留着一座中央峰，这些特征在重力场上都会跟着显出来。反观背面那几片低地，外缘的隆起是有，里面也确实平坦，可下面的月壳既没有变薄，中央也没有峰起。把边界和重力两件事摆到一块儿，能得出的读法就只有一种，虽然听上去不太寻常：那几片低地的底部原本和四周高地的顶部在同一个平面上，四周之所以高出一圈，是因为有人在边缘垒起过东西，中间原先并没有被撞击挖空。这是笔者手里这些数据能贴得最近的一种解释，也是眼下最难验证的一种。

\section{静海七号}

前面量过的所有东西全在月球背面，靠的都是轨道仪器传回来的照片、反照率和重力场数据，笔者本人没有去过月面，一块实地的石头也没有摸过。想再往下推，就必须先回到地面上拿得住的实物。

这类实物据说不止一次被带回地球。当年大概也算不上秘密，报纸登过消息，样品在几座城市巡展过，国会甚至还为此开过听证会。笔者把这段经过抄进稿子时自己都觉得有点恍惚，像是在转述别人记岔了的旧闻，因为麻烦出在另一头：到今天为止，还记得那些巡展和听证会的人，一个也找不到了。至于为什么谁都不记得，不在本书要讨论的范围之内，此处只记下这件摆在眼前的事实：实物一直都在，从来没有被藏起来过。\footnote{关于当年公开记录与大众记忆何以同时空白、文献何以无人追问，参见索科洛夫《记忆消除场的文献学后果》。}

留在手边的只有三只木箱，里面装着三十七片薄片和一份清单。箱子上的编目卡一律写着「远东出土文书」，结论只有八个字：来源不明，暂不归类。箱里另夹着一张卡片，写明这批物件出在月球正面 \texttt{MARE-TRANQ-07} 第 4 号居住环西侧储物间。不过除了这张卡片之外，再也找不到第二份独立记录，也查不出这批箱子究竟是哪一年、由哪一支队伍取回来的。那处设施后来有了一个通用编号，叫作静海七号，位于静海（Mare Tranquillitatis）东缘。\footnote{三十七片残简的释文与考释见《月之都——〈广寒宫残简〉辑佚与考释》，该书已经把此处的形制、净秽、岁计与撤离讲过一遍，这里不重复，只借用同箱的那份清单，编号 \texttt{ФВ-1963-17}。}

\begin{table}[htbp]
  \centering
  \caption{静海七号的尺寸}
  \begin{tabular}{@{}lll@{}}
    \toprule
    & 外重 & 内重\\
    \midrule
    形制 & 屏蔽环 & 垣 \\
    直径 & 约 3.5 公里 & 约 1.2 公里\\
    周长 & — & 九里（秦汉一里 415.8 米，合 3.74 公里）\\
    高 & — & 十二丈（合 27.7 米）\\
    \bottomrule
  \end{tabular}
  \tabsource{据 \texttt{ФВ-1963-17} 移交清单第一页。清单只记形制与尺寸，不记方位。}
\end{table}

表里的九里折合三点七四公里，与一点二公里的直径恰好能对上，前提是按秦汉一里四百一十五点八米来换算。这个换算方法倒不是笔者硬套的，因为清单原文自己写了「里」字，想折算成现代米制，手头也只有这一套尺度能用。至于这份清单自己靠不靠得住，箱里的物件还能互相核对：残简里反复出现的门宽和环距都与清单吻合，薄片的长度也跟清单录下的构件尺寸对得上号。两头都是实物，不是后人抄本，彼此咬合得住。话虽如此，这种互校也只能证明这箱东西内部自圆其说，证明不了卡片上写的那个地点；三十七片薄片、一份清单和一张卡片同出一只木箱，它们互相吻合，顶多说明整箱材料出自同一批工场，还证明不了那个储物间果真在月面上存在过。

现在把静海七号和背面那几片低地并排放在一起。围住背面低地的环壁以千公里计，而静海七号的外重屏蔽环直径只有三点五公里，两者差着三个数量级。面对这么大的尺度差距，大概只有两种解释：一种是认为背面低地根本属于另一类地质或工程现象，前面量出来的齐整边界另有来历；另一种是认为背面低地和正面属于同一套形制，静海七号只是其中很小的一角。档案里有一句话稍微偏向后一种读法：清单专门注明「垣外为空」，意思是屏蔽环之内、内垣之外的大片空地上什么建筑也没有画。建筑图留白本也寻常，可是一座直径只有三点五公里的设施，原本没必要在垣外特意空出那样宽一圈地皮。笔者暂且采用后一种读法，依据也就这一条，算作旁证还行，远远算不上铁证。

通章算下来，这一章实际只证明了两件事：月球背面那几片低地的边界带有人工划定的迹象，而正面编号静海七号的那处设施直径只有三点五公里。至于「背面低地里正是一座城」，这一章其实一个字也没有证明。从一处量得出尺寸的小遗址，跨到另一处没人实地走过的大低地，中间这一步是笔者自己猜出来的。笔者不觉得这一步毫无根据，只是它还没有经过实地检验。如果日后真有人在月球背面踩到一段修筑起来的环壁，这一推论就算立住了；在此之前，它只是一份前后还能接上的假说。至于这城里究竟住过什么人，这一章更无从谈起，留到第二章再接着往下说。
